\section{Источники погрешности в математическом моделировании}
\label{sec:error-sources-in-modeling}

\subsection{Погрешность математической модели}

Математическая модель компромисс между бесконечной сложностью изучаемого явления и желаемой простотой его описания. Вследствие этого решение, полученное численно с помощью постренной модели может отичаться от точного решения из-за погрешностей.

Источники погрешности:
\begin{enumerate}
	\item Математическая модель (приближенное описание реального
	процесса) + исходные данные (результат
	эксперимерта/решение вспомогательной задачи) $\rightarrow$
	погрешность задачи $\delta_1y$ (неустранимая погрешность, зависит
	от степени адекватности модели реальному процессу);
	\item Численный метод (приближенный) $\rightarrow$ погрешность метода $\delta_2y$
	\item Округление $\rightarrow$вычислительная погрешность $\delta_3y$
	(определяется характеристиками ЭВМ)
\end{enumerate}

Полная погрешность результата решения задачи:

$$ \delta y = y - y^* = \delta_1 y + \delta_2 y + \delta_3 y $$

где \( y \) и \( y^* \) - точное и численное решение соответственно.

$$ \delta_1 y > \delta_2 y > \delta_3 y $$

\subsection{Абсолютная и относительная погрешности}

Пусть \( y \) — точное (неизвестное) значение некоторой величины, \( y^* \) — приближенное (известное) значение той же величины.  
Абсолютная погрешность приближенного значения \( y^* \)

$$ \Delta(y^*) = |y - y^*| $$

Относительная погрешность (не зависит от масштаба величины, единиц измерения)

$$ \delta(y^*) = \frac{|y - y^*|}{|y|} $$

Значение \( y \) неизвестно \(\Rightarrow\) получение оценок погрешностей, или верхних границ абсолютной и относительной погрешностей

$$
\Delta(y^*) \leq \overline{\Delta}(y^*), \quad \delta(y^*) \leq \overline{\delta}(y^*) $$

\subsection{Погрешности арифметических операций}

Пусть \(a^*\) и \(b^*\) - приближенные значения чисел \(a\) и \(b\).

\textbf{Утверждение 1.}

$$ \Delta(a^* \pm b^*) \leq \Delta(a^*) + \Delta(b^*) $$

\textbf{Утверждение 2.}

Пусть \( a \) и \( b \) - ненулевые числа одного знака. Тогда

$$ \delta(a^* + b^*) \leq \delta_{max}, \quad \delta(a^* - b^*) \leq \frac{|a + b|}{|a - b|} \delta_{max},$$

где \( \delta_{max} = \max(\delta(a^*), \delta(b^*)) \).

\textbf{Утверждение 3.}

$$ \delta(a^*b^*) \leq \delta(a^*) + \delta(b^*) + \delta(a^*) \delta(b^*) \approx \delta(a^*) + \delta(b^*) $$

$$ \delta\left(\frac{a^*}{b^*}\right) \leq \frac{\delta(a^*) + \delta(b^*)}{1 - \delta(b^*)} \approx \delta(a^*) + \delta(b^*) $$

\subsection{Подходы к учёту погрешностей}

\begin{enumerate}
	\item Аналитический (точное оценивание погрешности черезе арифм. операции)
		Слишком громоздкие вычисления.
	\item Статистический (правило Чеботарева)
		$$\Delta(x) \approx \frac{1}{n} \sqrt{3n} \Delta(x_i) = \sqrt{\frac{3}{n}} \Delta(x_i) \xrightarrow{n \to \infty} 0.$$
		Может приводить к увеличению точности.
\end{enumerate}

\subsection{Представление вещественных чисел}

Представление с плавающей точкой (floating point):

$$ x = \pm (\gamma_1 2^{-1} + \gamma_2 2^{-2} + \dots + \gamma_t 2^{-t}) 2^p, $$

где \(\gamma_1, \dots, \gamma_t\) - двоичные цифры, причем \(\gamma_1 = 1\).

\begin{enumerate}
	\item Множество компьютерных чисел float дискретно. Все остальные числа \( x \) имеют приближенное представление $ x^* = fl(x) $  
	с ошибкой округления. Относительная погрешность представления равна  
	
	$$ 	\overline{\delta}(x^*) = 2^{1-t} $$
	\item Диапазон компьютерных чисел ограничен.
	\item На машинной числовой оси числа расположены неравномерно.
	\item Машинный эпсилон, $\varepsilon_M$ расстояние между единицей и
	ближайшим следующим за ней числом.
	
	Машинный эпсилон $\varepsilon_M$ служит мерой относительной точности
	представления вещественных чисел.
\end{enumerate}

\subsection{Что важно сказать на экзамене}

Привести 3 основных источника погрешностей, сказать про абсолютную/относительную погрешности. Сказать про представление вещественных чисел. Уметь пояснить, зачем вообще нам погрешности нужны.

\newpage